\clearpage
\newpage

\section{Justification for LLM-as-Judge Evaluation}\label{appendix:llm-as-judge}

Our decision to utilize an LLM-based evaluation protocol, specifically using Qwen2.5-72B-Instruct~\citep{qwen2.5}, instead of the strict Exact Match (EM) scripts provided by certain benchmarks, e.g., GAIA~\citep{mialon2023gaia}, is substantiated by three key factors: the open-ended nature of the agent's outputs, consistency with recent literature, and the empirical reliability of the judge model.

\textbf{Mismatch between Exact Match and Open-ended Generation:} Official evaluation scripts for benchmarks like GAIA often rely on strict string matching or heuristic rules. These metrics typically assume the model is prompted to output a constrained format, such as a single number or a short phrase. However, the primary goal of ReSum is to enable unbounded exploration and self-evolution, which naturally results in agents generating detailed, reasoned responses. Applying strict EM to these outputs leads to a high rate of \textit{false negatives}, where the agent's answer is semantically correct but rejected due to formatting discrepancies. LLM-as-a-Judge provides the necessary semantic flexibility to assess correctness in such open-ended settings.

\textbf{Alignment with Existing Baselines:} Recent leading works in web agents, including ASearcher~\citep{gao2025asearcher}, ARPO~\citep{dong2025arpo}, and WebSailor~\citep{li2025websailor}, have shifted towards LLM-based evaluation to handle the complexity of open-ended web tasks. Since our experiments aim to benchmark ReSum against these open-source agents, adopting the same evaluation protocol is essential to ensure fair and direct comparability.

\textbf{Reliability of Qwen2.5-72B-Instruct as a Judge:} To ensure that our choice of judge model does not introduce bias, we conducted a comparative analysis of four candidate models: Qwen2.5-72B-Instruct, GPT-4o-Mini, Qwen3-235B, and Gemini-2.5-Flash. We evaluated agent predictions on the BrowseComp-zh dataset using all four judges. The results demonstrated a high degree of consensus, with score variance across models being less than 0.3\%. Specifically, the scores were 12.8\% (Qwen2.5), 12.5\% (GPT-4o-Mini), 12.8\% (Qwen3), and 12.5\% (Gemini). Discrepancies were primarily attributed to cross-lingual alignment challenges, e.g., matching an English ground truth title to a Chinese translation, where Qwen2.5 demonstrated superior capability in identifying semantic equivalence. Consequently, Qwen2.5-72B-Instruct was selected for its reliability, high agreement with larger proprietary models, and cost-effectiveness.
